\section{模型评价与结论}
\subsection{模型优点}
本文以同一能量平衡和SOC递推贯穿四问，新增机制只改变信息集和费用项，便于归因比较。预测按日期严格前推，避免随机划分造成时间泄漏；优化前验证可行性，优化后独立回代1,337次并保存迭代记录。分位裕度将紧急购电代价转化为可解释的保守净负荷，计算规模适合全年多次滚动求解。
\subsection{模型局限}
分位鲁棒LP是完整场景风险模型的近似，未显式描述时序相关场景或输出CVaR。储能互斥依赖正电价、效率损耗和极小吞吐正则项，而非二进制变量；虽然实际检查未出现同时充放电，但在允许负电价或售电时需改为混合整数模型。滚动调整仅比较两个端点策略，尚不能确定最优发布频率。
\subsection{结论与推广}
确定性条件下，储能低充高放可使典型日费用下降26.90\%；不确定条件下，历史滚动预测与80\%残差分位裕度使334日费用下降22.96\%。多时点更新能减少紧急购电，却可能因调整罚金提高总费用；波动电价预测在问题4-2中仅带来约0.05\%的费用改善。因此实际运营应把“是否调整”作为带成本的决策，而不是收到新预报后机械重算。该框架可推广至园区综合能源、充电站和含可再生能源的企业微网，只需替换设备约束、结算规则和预测模块。
